\newpage
\chapter{PENDAHULUAN} \label{Bab I}

\section{Latar Belakang} \label{I.Latar Belakang}

Ketidakseimbangan distribusi kelas pada dataset menjadi salah satu permasalahan yang paling sering dihadapi dalam penerapan \textit{machine learning} untuk klasifikasi. Kondisi ini terjadi ketika jumlah sampel pada satu kelas jauh lebih banyak dibandingkan kelas lainnya, sehingga kelas dengan jumlah sampel sedikit disebut kelas minoritas dan kelas dengan jumlah banyak disebut kelas mayoritas \cite{He2009Imbalanced}. Permasalahan data tidak seimbang banyak ditemukan pada berbagai domain nyata seperti deteksi penipuan kartu kredit, diagnosis penyakit langka, pemantauan sistem industri, dan analisis sentimen \cite{Krawczyk2016Imbalanced}. Algoritma klasifikasi yang dilatih pada data tidak seimbang cenderung memprioritaskan kelas mayoritas dan mengabaikan kelas minoritas, padahal kelas minoritas seringkali menjadi target utama analisis \cite{Hou2025XGBoost}. Model yang dihasilkan mungkin memperoleh nilai akurasi tinggi secara keseluruhan, namun gagal mendeteksi sampel kelas minoritas yang seharusnya teridentifikasi \cite{VanFC2024Enhancing}.

Krawczyk (2016) mengategorikan pendekatan untuk menangani data tidak seimbang ke dalam tiga kelompok utama, yaitu metode tingkat data (\textit{data-level}), metode tingkat algoritma (\textit{algorithm-level}), dan metode hibrida \cite{Krawczyk2016Imbalanced}. Metode tingkat data bekerja dengan memodifikasi distribusi dataset pelatihan melalui penambahan sampel kelas minoritas (\textit{oversampling}) atau pengurangan sampel kelas mayoritas (\textit{undersampling}). Metode tingkat algoritma memodifikasi algoritma klasifikasi secara langsung agar tidak bias terhadap kelas mayoritas, misalnya melalui pendekatan \textit{cost-sensitive learning} \cite{Krawczyk2016Imbalanced}. Metode hibrida menggabungkan pendekatan \textit{data-level} dengan teknik \textit{ensemble} untuk menghasilkan model yang lebih \textit{robust} dan efektif dalam menangani distribusi kelas yang timpang \cite{Krawczyk2016Imbalanced}. Di antara ketiga pendekatan tersebut, metode tingkat data menjadi yang paling banyak diteliti karena bersifat independen terhadap algoritma klasifikasi dan dapat diterapkan tanpa mengubah proses pelatihan model \cite{He2009Imbalanced}.

Pada kategori \textit{data-level}, teknik \textit{oversampling} dan \textit{undersampling} merupakan dua pendekatan yang paling banyak digunakan. Chawla \textit{et al.} (2002) memperkenalkan SMOTE yang menghasilkan sampel sintetis kelas minoritas berdasarkan interpolasi antar tetangga terdekat, sehingga mampu meningkatkan kemampuan model dalam mengenali kelas minoritas \cite{Chawla2002SMOTE}. Bunkhumpornpat \textit{et al.} (2009) mengembangkan Safe-Level SMOTE yang memperhitungkan tingkat keamanan posisi sampel sintetis berdasarkan komposisi tetangga terdekatnya, sehingga sampel baru tidak terbentuk di wilayah yang didominasi kelas mayoritas \cite{Bunkhumpornpat2009SafeLevel}. Masing-masing teknik memiliki kelemahan, dimana \textit{oversampling} berpotensi menimbulkan \textit{overfitting} akibat data sintetis yang redundan, sementara \textit{undersampling} berisiko menghilangkan informasi penting dari kelas mayoritas \cite{Eroglu2024HOUM}.

Untuk mengatasi kelemahan penggunaan teknik tunggal, Eroglu dan Pir (2024) mengusulkan HOUM (\textit{Hybrid Oversampling and Undersampling Method}) yang menggabungkan Safe-Level SMOTE untuk \textit{oversampling} dan SVM untuk \textit{undersampling} secara iteratif hingga dataset mencapai keseimbangan \cite{Eroglu2024HOUM}. HOUM memanfaatkan SVM untuk mengidentifikasi dan menghapus data kelas mayoritas yang berada jauh dari \textit{decision boundary}, sehingga data yang tersisa lebih informatif bagi proses klasifikasi. Pengujian pada delapan dataset menunjukkan bahwa HOUM mampu mengungguli metode pembanding seperti SMOTE, SMOTEENN, SMOTETomek, W-SIMO, dan RusAda \cite{Eroglu2024HOUM}. Namun, penelitian tersebut hanya menggunakan satu konfigurasi parameter SVM dan satu teknik \textit{oversampling} (Safe-Level SMOTE) tanpa mengeksplorasi pengaruh variasi parameter terhadap hasil klasifikasi.

Performa SVM sangat dipengaruhi oleh pemilihan parameter regularisasi $C$ dan jenis kernel yang digunakan \cite{Cervantes2020SVM}. Parameter $C$ mengontrol keseimbangan antara kompleksitas model dan toleransi terhadap kesalahan klasifikasi, sedangkan jenis kernel menentukan cara SVM memetakan data ke dimensi yang lebih tinggi. Variasi konfigurasi parameter tersebut berpengaruh langsung terhadap pembentukan \textit{decision boundary} yang digunakan HOUM untuk proses \textit{undersampling}. Selain konfigurasi SVM, pemilihan teknik \textit{oversampling} berpotensi memengaruhi kualitas data hasil penyeimbangan. He \textit{et al.} (2008) memperkenalkan ADASYN (\textit{Adaptive Synthetic Sampling}) yang menghasilkan lebih banyak sampel sintetis pada area yang sulit diklasifikasikan berdasarkan distribusi bobot adaptif \cite{He2008ADASYN}. Pendekatan ADASYN berbeda dengan Safe-Level SMOTE karena fokus pada sampel minoritas yang sulit dipelajari, sehingga batas keputusan model dapat bergeser secara adaptif ke arah yang tepat \cite{He2008ADASYN}. Perbandingan antara kedua teknik ini sebagai komponen \textit{oversampling} pada HOUM belum pernah dilakukan pada penelitian sebelumnya.

Berdasarkan uraian tersebut, penelitian ini bertujuan menganalisis pengaruh variasi konfigurasi pada metode HOUM terhadap performa klasifikasi data tidak seimbang. Pada komponen \textit{undersampling}, dilakukan variasi parameter $C$ pada SVM dengan enam nilai (0,5; 1; 1,5; 2; 2,5; 3) serta perbandingan tiga jenis kernel yaitu Linear, RBF, dan Polynomial. Pada komponen \textit{oversampling}, dilakukan perbandingan antara HOUM yang menggunakan Safe-Level SMOTE dengan HOUM yang menggunakan ADASYN. Performa klasifikasi dievaluasi menggunakan metrik ROC-AUC, G-mean, \textit{Recall}, dan \textit{Precision} yang sesuai untuk menilai kemampuan model dalam menangani data tidak seimbang. Hasil penelitian ini diharapkan dapat memberikan rekomendasi konfigurasi HOUM yang optimal untuk berbagai skenario ketidakseimbangan data.


\section{Rumusan Masalah} \label{I.Rumusan Masalah}

Berdasarkan latar belakang, permasalahan penelitian dirumuskan sebagai berikut: \par

\begin{enumerate}[noitemsep]
	\item Bagaimana pengaruh penggunaan teknik \textit{oversampling} yang berbeda pada metode HOUM terhadap kinerja klasifikasi data tidak seimbang?
	\item Bagaimana pengaruh variasi nilai parameter regularisasi (C) dan jenis kernel \textit{Support Vector Machine} terhadap kinerja metode HOUM dalam menangani data klasifikasi tidak seimbang?
\end{enumerate}


\section{Tujuan Penelitian} \label{I.Tujuan}
Berdasarkan rumusan masalah, tujuan dari penelitian adalah sebagai berikut: \par

\begin{enumerate}[noitemsep]
	\item Menganalisis pengaruh penggunaan teknik \textit{oversampling} yang berbeda pada metode HOUM terhadap kinerja klasifikasi data tidak seimbang. 
	\item Menganalisis pengaruh variasi parameter regularisasi (C) dan jenis kernel \textit{Support Vector Machine} terhadap kinerja metode HOUM dalam menangani data klasifikasi tidak seimbang.
\end{enumerate}


\section{Batasan Masalah} \label{I.Batasan}
Adapun batasan masalah dari penelitian  agar sesuai dengan yang diharapkan adalah sebagai berikut: \par

\begin{enumerate}[noitemsep]
    \item Penelitian difokuskan pada data klasifikasi tidak seimbang (\textit{imbalanced classification}) dan tidak mencakup permasalahan regresi.
    \item Teknik \textit{oversampling} yang dianalisis pada metode HOUM dibatasi pada Safe-Level SMOTE dan ADASYN.
    \item Analisis pengaruh konfigurasi SVM dibatasi pada pemilihan jenis kernel, yaitu Linear, \textit{Radial Basis Function} (RBF), dan Polynomial, serta variasi nilai parameter regularisasi (C).
    \item Nilai parameter regularisasi (C) yang digunakan dalam penelitian ini dibatasi pada nilai 0,5; 1; 1,5; 2; 2,5; dan 3.
    \item Evaluasi performa metode HOUM dibatasi menggunakan metrik ROC-AUC, G-Mean, \textit{Recall}, dan \textit{Precision}.
\end{enumerate}


\section{Manfaat Penelitian} \label{I.Manfaat}
Adapun manfaat yang diperoleh dari hasil penelitian adalah sebagai berikut: \par

\begin{enumerate}[noitemsep]
    \item Memberikan rekomendasi mengenai teknik \textit{oversampling} yang paling efektif (Safe-Level SMOTE vs ADASYN) untuk diintegrasikan dalam metode HOUM saat menangani data berimbang.
    \item Menjadi acuan dalam menentukan nilai parameter regularisasi ($C$) yang optimal pada SVM untuk meningkatkan kemampuan model HOUM dalam mengklasifikasikan kelas minoritas.
    \item Memberikan panduan empiris dalam memilih jenis kernel SVM yang paling sesuai untuk memaksimalkan performa \textit{undersampling} metode HOUM pada karakteristik data yang berbeda.
\end{enumerate}


\section{Sistematika Penulisan} \label{I.Sistematika}
Penulisan tugas akhir disusun secara sistematis ke dalam lima bab yang saling berkaitan. Berikut adalah uraian  mengenai isi masing-masing bab. \par

\noindent\textbf{Bab I Pendahuluan}

Pada bab ini menyajikan gambaran umum penelitian yang meliputi latar belakang permasalahan, rumusan masalah, tujuan penelitian, batasan masalah, manfaat penelitian, serta sistematika penulisan laporan tugas akhir. \par

\noindent\textbf{Bab II Tinjauan Pustaka}

Pada bab ini menguraikan kajian literatur yang berkaitan dengan topik penelitian, mencakup tinjauan terhadap penelitian-penelitian sebelumnya serta landasan teori yang mendasari metode dan konsep yang digunakan dalam penelitian ini. \par

\noindent\textbf{Bab III Metodologi Penelitian}

Pada bab ini menjelaskan secara rinci metodologi yang diterapkan dalam penelitian, meliputi tahapan alur kerja, perangkat dan dataset yang digunakan, metode yang diimplementasikan, serta rancangan skenario pengujian yang akan dilakukan. \par

\noindent\textbf{Bab IV Hasil dan Pembahasan}

Pada bab ini memaparkan hasil penelitian yang diperoleh dari implementasi metode, beserta analisis dan pembahasan terhadap hasil pengujian. \par

\noindent\textbf{Bab V Penutup}

Pada bab ini berisi kesimpulan yang merangkum hasil penelitian berdasarkan tujuan penelitian, serta saran-saran untuk pengembangan dan penyempurnaan pada penelitian selanjutnya.